\section{Kinetic energy, scalar relativity, and spin-orbit coupling}
\label{Sec:rel}

For elements beyond approximately $Z$=30, relativistic effects near
the nucleus cannot be neglected in an all-electron treatment---both
for core, and for valence electrons. For the purposes of ``everyday'' matter,
the full theory is given by Dirac's four-component Equation, but in the
``practice'' of materials physics and chemistry, we still tend to think in
terms of Schr\"odinger-like objects. A true four-component treatment is not
implemented in FHI-aims, but the following standard levels of approximation
are available: 
\begin{itemize}
  \item Non-relativistic kinetic energy (one or two collinear \keyword{spin}
    components of the Kohn-Sham orbitals)
  \item Scalar-relativistic kinetic energy expression (one or two collinear
    spin components). 
  \item Perturbative spin-orbit coupling, a single correction step to the
    Kohn-Sham eigenvalues based on the Kohn-Sham orbitals from a
    non-relativistic or scalar-relativistic s.c.f. cycle. \emph{Perturbative
      spin-orbit coupling in FHI-aims is a new feature. Importantly,
      changes to the total  
      energy beyond the sum-of-eigenvalues are not included, and total energy 
      gradients (forces) are also unavailable.}
\end{itemize}

\subsection*{Scalar relativity}

While the non-relativistic level of theory is exactly
defined and will be the same in any first-principles implementation (at a
complete basis set, all-electron level anyway), there are many different
versions of scalar-relativistic approximations which can yield
considerably different \emph{total} energies for different systems. Their
unifying feature is that any two scalar-relativistic methods should still
yield the same energy \emph{differences} for properties that concern
valence electrons: Binding energies, valence eigenvalues, etc. 

As outlined in Ref. \cite{Blum08}, FHI-aims offers two reliable, effective
one-component scalar relativistic treatments (by ``effective one-component'',
we mean that no coupling exists between the two spin channels of a calculation
with collinear spin): ``atomic ZORA'' and ``scaled ZORA''. (ZORA: ``zero-order
regular approximation''.) The potential energy surfaces and valence
eigenvalues obtained from either one are sound, with total energy
\emph{differences} (e.g., between different structures) of atomic ZORA
matching benchmark results obtained by other scalar-relativistic
methods somewhat more closely in the ``Delta Project'',
https://molmod.ugent.be/deltacodesdft .   

By construction, scaled ZORA (absolute) \emph{total} energies and
\emph{core} eigenvalues are closer to their counterparts from
other scalar-relativistic methods than atomic ZORA. 

However, in practical calculations, atomic ZORA appears to be
preferable: 
\begin{itemize}
  \item The computational overhead for scaled ZORA is significant
  \item Calculating energy gradients (forces) is much more
    involved than for atomic ZORA, which is why gradients are only
    implemented for atomic ZORA
  \item It appears (see above) that atomic ZORA based energy
    differences are slightly more reliable than scaled ZORA energy
    differences. 
\end{itemize}

In practice, any valence electron related results from atomic ZORA,
scaled ZORA, or other scalar relativistic methods in the literature
should be expected to agree with one another up to benchmark accuracy
level if all \emph{other} computational tasks are numerically well
converged. Only the broad comparison of materials and codes in the
``Delta Project'' (https://molmod.ugent.be/deltacodesdft) helped us
distinguish what is otherwise a very subtle difference at best. 

As illustrated in more detail in Sec. \ref{Sec:rel-example}, we
presently recommend the following strategy for the correct handling
of scalar relativity for arbitrary structures in FHI-aims:
\begin{enumerate}
  \item For energy differences, geometries, valence eigenvalues etc.,
    simply rely on the \emph{atomic} zero-order regular approximation
    approximation (ZORA). This has performed very well for us in a
    broad array of benchmarks.
  \item If you wish for a ``second opinion'' or, for some reason, need
    a more credible (still not accurate!) handling of core
    eigenvalues, perform single-point calculations using the ``scaled
    ZORA'' \cite{Lenthe94} approximation. 
\end{enumerate}
Similarly, any other tasks involving forces should be performed using
atomic ZORA, remembering that only energy \emph{differences} between
different structures are meaningful in this case.

\subsection*{Spin-Orbit Coupling}

FHI-aims currently implements a treatment of the spin-orbit coupling
operator which adds spin-orbit coupling (SOC) corrections to the Kohn-Sham
eigenvalues, band structures, and densities of states in a
single evaluation after a scalar-relativistic s.c.f. cycle without
spin-orbit coupling is complete. In other words, this is a
non-selfconsistent evaluation of spin-orbit coupling. The keyword 

  \keyword{include\_spin\_orbit} \option{non\_self\_consistent}

should be all that is needed for most production applications. This keyword
should be \emph{added} to a standard \texttt{control.in} file.
The level of scalar relativity should still be specified as usual, for
example, 

  \keyword{relativistic} \option{atomic\_zora} \option{scalar}

to invoke the atomic ZORA kinetic energy operator.

Applying the SOC operator in this perturbative
(non-selfconsistent) way is quantitatively accurate (to a few 0.01 eV
for valence band structures) for elements below Xe ($Z$=54) when
combined with atomic ZORA.

\emph{For heavy
  elements (approximately Au and beyond) this level of theory is only
  qualitatively accurate. It captures the majority of the SOC effect,
  but quantitative deviations above 0.1 eV for band structures must be
  expected. Similarly, the corrections for any core levels would
  require one to go beyond non-selfconsistent SOC.}

Since FHI-aims, by default, only treats a finite number of bands in a
given calculation, increasing the number of bands using the keyword
\keyword{empty\_states} may be necessary to capture correctly the
behavior of high-lying conduction bands.

\subsubsection{Some Background}

Spin-orbit coupling is a simple consequence of transforming Dirac's
equation to a (two-component) Schr\"odinger-like form. This leads to
an approximate Hamiltonian of the form
\begin{equation}
  \hat{H} = \hat{t}_{SR} + \hat{v} + \hat{v}_{SOC} .
\end{equation}
$\hat{t}_{SR}$ is the usual scalar relativistic kinetic energy
operator (e.g., atomic ZORA). $\hat{v}$ is the local or non-local
potential, and $\hat{v}_{SOC}$ is the spin-orbit coupling operator,
\begin{equation}
  v_{SOC}=\frac{i}{4c^{2}}\bm{p} \tilde{v} \times \bm{p} \cdot
  \bm{\sigma} .
\end{equation}
This operator also contains a potential $\tilde{v}$, and in principle, for
Kohn-Sham theory (local potential), $\tilde{v}$ could be equal to
$v(r)$. In practice, however, the combination of potential and
gradients in $v_{SOC}$ means that SOC arises essentially only very
near the nucleus, where the potential is deep and dominated by the
nuclear potential. In FHI-aims, we therefore only consider the sum of
the nuclear and the electrostatic potentials, but not the Kohn-Sham
exchange-correlation potential, for $\tilde{v}$ when constructing $v_{SOC}$.

We then treat spin-orbit coupling $\bolds\times\boldp$ perturbatively
, i.e., as a single-shot postprocessing step after an otherwise
converged scalar-relativistic self-consistent field cycle. 
By default, we can simply include all Kohn-Sham states in the
evaluation of $v_{SOC}$ above and invert the resulting Hamiltonian
once to obtain its eigenvalues and eigenvectors. This amounts to
standard first-order perturbation theory and is the level of theory
that we have found to be sufficient for accurate results.

It is possible to go beyond the first-order level by so-called
quasi-degenerate (qd) perturbation theory in a second-variational
approach. Some of this is implemented in FHI-aims and summarized below
for completeness. However, we have not found any significant
improvements compared to the first-order level. At present, it does
not appear necessary to play with any of this except for development
or research purposes. 

\emph{Quasi-degenerate perturbation theory} is based on L\"owdin partitioning
\cite{Loewdin51}. The spin-orbit interaction is seen as a perturbation to the 
scalar-relativistic Kohn-Sham-Hamiltonian $\hat{h}^\text{sr}$:
\begin{equation}
\label{eqn:soc_ham}
\hat{h}=\hat{h}^\text{sr}+\hat{h}^\text{soc} \, .
\end{equation}
The scalar-relativistic set of eigenstates $\left\{\psi^\text{sr}\right\}$ is
then divided into several subsets $\left\{\psi^\text{sr}\right\}_k$ in the
following way. Two scalar-relativistic states $i$ and $j$ belong to the same
subset if   
\begin{enumerate}
 \item their eigenvalues are close OR
 \item their matrix element $\left<\psi_i^\text{sr}\left|\hat{h}^\text{SOC}\right|\psi_j^\text{sr}\right>$ 
	is  comparable to the eigenvalue difference
\end{enumerate}
which means that any two \emph{different} subsets are weakly interacting in
terms of $h^\text{SOC}$ and separated in energy. For each subset the
Hamiltonian \ref{eqn:soc_ham} is then 
``block-diagonalized'' to a form $\tilde{h}$ such that all matrix elements 
$\left<\psi_{\mathrm{in\;subset}}\left|\tilde{h}\right|\psi_{\mathrm{not\;in\;subset}}\right>$  
vanish up to a specified order in $h_{\mathrm{SOC}}$. For details see reference
\cite{Winkler:SOC}. 


\newpage

\subsection*{Tags for general section of \texttt{control.in}:}

\keydefinition{include\_spin\_orbit}{control.in}
{
  \noindent
  Usage: \keyword{include\_spin\_orbit} \option{method} \\[1.0ex]
  Purpose: Calculates spin-orbit coupling corrected eigenvalues. \\[1.0ex] 
  \option{method} The method for including spin-orbit coupling.  At present, only \option{type} 
  \option{non\_self\_consistent} is suitable for production-level calculations. \\[1.0ex] 
}

\emph{Note: While FHI-aims also prints out a corrected total-energy
  expression based on the SOC-corrected eigenvalues, please do not
  trust this total energy expression for anything. It is experimental.} 

Spin-orbit corrections can be added as a perturbative post-processed
step after a non-relativistic, atomic ZORA or scaled ZORA
scalar-relativistic calculation. Usually, just writing

  \keyword{include\_spin\_orbit} \option{non\_self\_consistent}

is sufficient to obtain corrected eigenvalues and (if requested) band
structures, densities of states, etc., in first-order perturbation
theory to a scalar-relativistic calculation.

FHI-aims will write the results to standard output after the SCF cycle
has converged and details to a separate output file
\option{SOC\_result.out}.  

Since perturbative spin-orbit coupling operates
on the space spanned by the calculated scalar-relativistic
eigenvectors, for accurate high-lying bands one must include
sufficiently many unoccupied states. If needed, this may be done by
increasing the keyword \keyword{empty\_states}. 

\emph{For research purposes, FHI-aims also includes some level of theory
beyond first-order perturbation theory. This is summarized below and
gives rise to some additional keywords. However, we have not seen a
large improvement from the approaches described below.}

FHI-aims also enables the use of (quasi-)degenerate perturbation
theory, for which there are three important parameters:
\option{matrix\_tolerance} and \option{ev\_tolerance} which  
determine the assignment of eigenvalues to degenerate subspaces, and \option{soc\_perturb\_order} which determines the 
order of (quasi-)degenerate perturbation theory to perform. 
 
\option{matrix\_tolerance} places eigenstates in the same subspace if \\ 
    $\left|\left<\psi_i^\text{sr}\left|\hat{h}^\text{soc}\right|\psi_j^\text{sr}\right>\right|\times
    \option{matrix\_tolerance}\ge \left|\epsilon_i-\epsilon_j\right|$ \\
\option{ev\_tolerance} [in eV] places eigenstates in the same subspace if 
    $\left|\epsilon_i-\epsilon_j\right| \le \option{ev\_tolerance}$

Standard first-order perturbation theory is a special case of the
quasi-degenerate perturbation formalism, achievable by considering
only first-order effects and setting a large enough value of
\option{ev\_tolerance} that all eigenstates are placed into a single
subspace, and we have chosen to make this the default setting. 

While we believe standard first-order perturbation theory gives
reliable numbers as compared to to quasi-degenerate perturbation
theory and highly recommend the user stay with it, nevertheless if the
user wishes to change these settings, they may do so.  The default
value for \option{matrix\_tolerance} is 0.1.  The default value for
\option{ev\_tolerance} is 1000000000 eV.  To change these, use keyword
\keyword{soc\_subspace}.  The default value for
\option{soc\_perturb\_order} is 1.  To change this, use keyword
\keyword{soc\_perturb\_order}. 


\emph{If the user wishes to play around with quasi-degenerate
  perturbation theory,   we have found in the past that
  \option{soc\_perturb\_order} 2 and   \option{ev\_tolerance} 0.00001
  eV are relatively safe starting values, and that
  \option{matrix\_tolerance} is the parameter that should adjusted,
  but we give the    user the ability to change all three.  The
  characteristic signature of poorly    chosen values are ``wiggles''
  in the plotted band structure.} 

\keydefinition{soc\_subspace}{control.in}
{
  \noindent
  Usage: \keyword{soc\_subspace} \option{parameter} \option{value} \\[1.0ex]
  Purpose: Used to change the tolerance parameters associated with subspace assignment for perturbative SOC. \\[1.0ex]
  \option{parameter} The parameter to change.  Options are \option{matrix\_tolerance} or \option{ev\_tolerance}. 
     \\[1.0ex]
  \option{value} The value assigned to the parameter chosen by \option{parameter}. \\[1.0ex]
}
Default:  \option{matrix\_tolerance} 0.1, \option{ev\_tolerance} 1000000000 eV.

This is an optional keyword.  See \keyword{include\_spin\_orbit} for more information.

\keydefinition{soc\_perturb\_order}{control.in}
{
  \noindent
  Usage: \keyword{soc\_perturb\_order} \option{order} \\[1.0ex]
  Purpose: Sets the order of quasi-degenerate perturbation theory to use for SOC calculations. \\[1.0ex] 
  \option{order} The order of quasi-degenerate perturbation theory to use, an integer from 0 to 3. \\[1.0ex] 
}
Default: 1.

This is an optional keyword.  See \keyword{include\_spin\_orbit} for more information.

\keydefinition{compute\_kinetic}{control.in}  
{
  \noindent
  Usage: \keyword{compute\_kinetic} \\[1.0ex]
  Purpose: \emph{Experimental} - for test purposes, allows to compute the kinetic energy via
    the product of the kinetic energy matrix and the density matrix
    \\[1.0ex]
}
This flag is presently kept for test purposes only (the electronic
kinetic energy is separately computed and printed for each scf
iteration anyway) but may be useful for some future modifications.

\keydefinition{override\_relativity}{control.in}
{
  \noindent
  Usage: \keyword{override\_relativity} \option{flag} \\[1.0ex]
  Purpose: If explicitly set, allows to override the stop enforced by
  the code when physically questionable \keyword{relativistic}
  settings are used.  \\[1.0ex]
  \option{flag} is a logical expression, either \texttt{.true.} or
    \texttt{.false.} Default: \texttt{.false.} \\
}
For example, this will allow you to run a physically incorrect
calculation of heavy elements (think Au) with Schr\"odinger's
expression for the kinetic energy, instead of a scalar relativistic
treatment. The results will be wrong, so this flag should only be set
for test purposes. When set, the code assumes that the user must know
what they are doing.

\keydefinition{relativistic}{control.in}
{
  \noindent
  Usage: \keyword{relativistic} \option{r-type} \option{s-type}
    [\option {threshold}] \\[1.0ex]
  Purpose: Specifies the level of relativistic treatment in the
    calculation. \\[1.0ex]
  \option{r-type} is a string, specifying the basic approximation
    made. \\
  \option{s-type} is a string, specifying whether a scalar treatment
    is desired (currently, only the \texttt{scalar} option is
    supported). \\ 
  \option{threshold} is a small positive real number, allowing to
    reduce some integration effort. \\
}
Detailed expressions for the scalar relativistic treatments available
here are given in Ref. \cite{Blum08}. We here only repeat the salient
options and expressions. Possible options for \option{r-type} are:
\begin{itemize}
  \item \texttt{none} : Non-relativistic kinetic energy. In this case,
        \option{s-type} and \option{threshold} need not be provided. 
  \item \texttt{atomic\_zora} : Atomic ZORA approximation as described
    in Ref. \cite{Blum08}. \option{threshold} need not be provided. 
    This is the currently recommended option for energy differences and valence 
    and unoccupied eigenvalues.
  \item \texttt{zora} The ZORA approximation is used throughout the
    s.c.f. cycle, followed by a ``scaled ZORA'' \cite{Lenthe94} post-processing step
    (rescaling of all eigenvalues). \textbf{WARNING:} Do not rely on
    intermediate, simple ZORA total energies, but only on the final,
    rescaled total energies instead! ZORA (unscaled) is not the same as 
    ``atomic ZORA'' and cannot be trusted.
\end{itemize}
Forces are only provided for \texttt{none} and \texttt{atomic\_zora}.

\textbf{Remember to never take energy differences between calculations
  performed with different ``relativistic'' settings.}

In case of \texttt{zora}, the \option{threshold} option is required;
it specifies the threshold value above which the difference between
the sum-of-free-atoms ZORA expression and that for the actual
potential during the s.c.f. cycle will be calculated. In areas of
shallow potentials, where both expressions are substantially similar,
this saves the extra integration effort associated with ZORA. For
(very!) safe settings, \option{threshold} may be set to 10$^{-12}$; in
our experience, also 10$^{-9}$ does not lead to any noticeable
accuracy loss. 

Default is \texttt{none} if all elements in the structure have $Z<$20
(no heavier than Ca). For all heavier elements, an explicit
\keyword{relativistic} setting is required. For 11$<Z<$20, the setting
\texttt{none} will be accepted, but a
warning will be issued. For $Z$>20, choosing \texttt{none} will cause
the code to stop with an error message in order to avoid accidental
calculations with incorrect relativity. If a non-relativistic
calculation is still desired, for example for test purposes, this
``stop'' can be disabled by setting the flag
\keyword{override\_relativity}---but use this only if you know what
you are doing.

